%\documentclass{cumcmthesis}
\documentclass[withoutpreface,bwprint]{cumcmthesis} %去掉封面与编号页，电子版提交的时候使用。
\usepackage{booktabs}
\usepackage{tabularx}
\usepackage{ragged2e} % 优化换行文本对齐

\usepackage[framemethod=TikZ]{mdframed}
\usepackage{url}   % 网页链接
\usepackage{subcaption} % 子标题
\title{快递公司送货策略数学建模论文}
\tihao{A}
\baominghao{4321}
\schoolname{广东工业大学大学}
\membera{ 王宁}
\memberb{ 吴美欣}
\memberc{杨青梅 }
\supervisor{金应华 }
\yearinput{2026}
\monthinput{08}
\dayinput{12}

\begin{document}

 \maketitle
 \begin{abstract}
针对快递企业 30 个末端配送点调度优化问题，本文建立带载重、工时双重硬约束的多行程车辆路径（MVRP）模型，采用两阶段分层启发式算法求解：先以最远点优先贪心规则(最远点优先贪心为本文自定义的启发式构造规则，并非 VRP 经典标准算法)划分满足载重限制的配送趟次，再通过 FFD 装箱法将多趟配送任务分配给业务员，配套组内 2-opt 最短路排序与 MOVE/SWAP 跨组调点局部搜索优化配送成本。
\begin{itemize}
    \item 问题一以最小化业务员数量为目标，约束单人每日工作不超过 6 小时、单车最大载重 25kg，程序求解得到共 8 个配送趟次，最少配置 5 名业务员，全程总行驶里程 478.0km；
    \item 采用载货 3 元 /(km・kg)、空载 2 元 /km 分段计价规则，以总薪酬最低为目标优化配送顺序与分组，通过 MOVE/SWAP 跨组调点局部搜索迭代优化后，最优方案共 9 趟配送，需 6 名业务员，配送总费用 13800.7 元；
    \item 开展工时参数灵敏度分析，将业务员单日工时上限由 6h 放宽至 8h，复用整套模型算法重新求解最小人力方案，配送趟次保持 8 趟不变，所需业务员由 5 人降至 4 人，总行驶里程维持 478.0km。
\end{itemize}

整套模型采用递进复用基础框架，仅调整目标函数或工时参数即可完成多场景计算；分层启发式算法计算高效，适配 30 点中小规模配送场景，但仅能输出较优可行解，可引入遗传算法、整数规划进一步逼近全局最优。



\keywords{MVRP车辆路径问题\quad  最远点优先贪心算法\quad   局部搜索\quad 载重分段成本\quad 配送调度优化\quad 参数灵敏度分析}
\end{abstract}

%目录  2019 明确不要目录，我觉得这个规定太好了
%\tableofcontents

%\newpage

\section{问题重述}
\subsection{问题背景}
快递企业每日从总部仓库统一配送快件至 30 个固定末端网点，配送路线为平行坐标轴的折线，采用曼哈顿距离核算行驶里程。
\par 配送存在两项核心运营约束：配送车辆单次载货上限 25kg、每名业务员单日存在最长工作时长限制；每个配送网点固定停留 10 分钟卸货。
\par 企业提出三层优化需求：一是在 6 小时工时限制下合理规划配送路线，尽可能减少业务员用工数量；二是基于载货、空载差异化计价规则优化配送顺序，降低整体薪酬成本；三是分析单人工作时长由 6 小时延长至 8 小时后，人力配置、配送里程、工作负荷的变化规律，为企业排班管理提供量化决策依据。30 个配送点快件总重量固定为 184.5kg，各网点坐标、单件快件重量已知。

\subsection{已知条件梳理}

\begin{itemize}
    \item 基础配送参数：配送原点（仓库）编号 0；共 30 个配送点，各点快件重量已知；单车最大载重 25kg；每个网点固定停留 10 分钟；业务员单日最大工作时长基础约束 6 小时；行驶距离采用曼哈顿距离。
    \item 问题一行驶规则：全程行驶速度 25km/h，路线需闭合（仓库出发，配送完成返回仓库），目标最小化业务员数量。
    \item 问题二计价与速度规则：载货行驶速度 20km/h，空载返程速度 30km/h；载货里程成本 3 元 /(km・kg)，空载返程里程成本 2 元 /km；目标最小化全部业务员总薪酬费用，载重、6 小时工时约束不变。
    \item 问题三调整规则：将单人每日工作时长放宽至 8 小时，重新求解最小人力配送方案，对比 6 小时、8 小时工时约束下业务员数量、总配送里程、单人工作负荷的变化，分析工时约束对配送调度的影响规律。
    \item 基础数据：30 个配送点总快件重量 184.5kg，坐标、单件重量固定。
\end{itemize}

\subsection{待求解问题}
（1）在载重≤25kg、单日工作≤6h 约束下，划分配送分组，确定业务员数量、每组闭合配送路线、各组总重量、总行驶里程、单组总工作时长；
\par（2）基于差异化载货 / 空载计价规则，先计算（1）分组方案总薪酬；再优化分组与配送顺序，得到费用最优配送方案，对比优化前后成本差异；
\par（3）将单日最大工作时长放宽至 8 小时，分别求解最小人力、最小成本两套调度方案，对比 6h、8h 约束下业务员数量、总里程、总薪酬的变化趋势（本文暂未完成求解，预留完整建模与求解框架）。

\section{问题分析}
\subsection{问题一分析}

1.	问题本质：本题业务员可以多次从仓库出发完成多条配送回路，属于多行程车辆路径 (MVRP)，区别于单车仅单条路线的经典 CVRP：带载重 + 工时双重硬约束的 MVRP 车辆路径优化问题，属于 NP 难组合优化；优化目标为最小启用业务员（车辆）数量，需同时满足重量上限与总工作时间限制。
\par2.	核心难点：分组满足载重约束时极易出现路线总工时超时，两类约束相互牵制，简单按距离聚类会产生大量无效分组。
\par3.	算法选型分析： 
 \begin{itemize}
  \item  简易方案：最远点优先贪心分组，逻辑简单、易编程实现，适配 30 点小规模场景，可快速输出可行解；
  \item  进阶方案：Clarke-Wright 节约算法，进一步缩减总行驶里程；
  \item  高精度方案：整数规划 / 遗传算法，求解速度慢，竞赛优先使用贪心算法作为基础求解工具。
\end{itemize}

本文选择最远点优先贪心分组 + 组内最近邻排序作为核心求解算法。

\subsection{问题二分析}

1.	问题本质：载重依赖型成本最优 VRP，区别于传统最小里程优化，载货成本与路线累计剩余重量线性相关，总行驶距离最短不等于配送薪酬最低。
\par2.	成本逻辑拆解：车辆行驶过程中，每经过一个配送点卸载货物，后续路段载货重量持续下降；重货后置会大幅抬高前段载货里程成本，因此需要在配送分组与访问顺序上协同优化，使高载重路段尽可能缩短，降低载货段的重量里程累积成本。
\par3.	优化思路： 
 \begin{itemize}
  \item  基线计算：直接复用问题一原始分组路线，分段核算载货、空载里程，计算总薪酬作为对比基准；
  \item  局部搜索优化：采用 MOVE（跨组移动配送点）、SWAP（跨组交换配送点）两类邻域操作，在不破坏载重、工时约束的前提下降低总费用；组内访问顺序由最近邻构造加 2-opt 优化按曼哈顿最短路确定。
 \end{itemize}

\subsection{问题三分析}
1.	模型逻辑：无需重构数学模型，仅修改工时约束参数（6h 调整为 8h），完整复用问题一的分趟、打包算法，属于约束参数的灵敏度实验。
\par2.	预期变化规律：单人允许工作时长提升，单名业务员单日可承载更多配送趟次，整体所需业务员数量减少；在本算例保持原有配送趟次分组不变的前提下，总行驶里程维持 478.0 km 不变，该结果为当前分组下的特定结果，不具备普适性，若工时放宽后重新优化配送点分组，每趟路径结构将发生改变，总行驶里程也会随之变化；代价为单人员单日配送负荷显著提升。
\par3.	对比分析维度：统一提取业务员总数、总行驶里程、单人均工作时长、车辆装载率四项指标，量化工时放宽带来的人力配置变化，为企业平衡用工成本与员工劳动强度提供依据。


\subsection{整体求解逻辑框架}
三问共用同一套基础数据、距离计算公式、载重约束，模型呈递进复用关系：

\begin{itemize}
  \item  问题一搭建基础约束模型，输出最小人力分组方案，为问题二提供基准对比解；
  \item  问题二复用全部载重、工时约束，仅更换优化目标为最小薪酬，新增分段成本核算逻辑；
  \item  问题三完全复用前两问模型，仅修改工时参数，完成两组参数下结果横向对比；
\end{itemize}
	
全文采用统一 Python 模块化代码实现计算，距离、工时、成本封装通用函数，保证三组方案指标具备横向可比性。

\section{模型假设与符号说明}
\subsection{模型假设}
1.	配送行驶路径严格采用曼哈顿距离，不考虑道路绕行、红绿灯、交通拥堵等额外耗时；
\par2.	每个配送点快件不可拆分，一个网点仅由一名业务员配送；
\par3.	所有业务员统一从仓库原点出发，完成配送后必须返回原点，路线为闭合回路；
\par4.	每个配送点固定停留 10 分钟卸货，停留时长与快件重量无关；
\par5.	所有配送任务静态已知，无临时新增、取消配送订单；
\par6.	车辆行驶匀速，载货、空载速度严格按照题目给定数值计算；
\par7.	不考虑业务员休息、车辆故障、油耗等额外隐性成本，仅计算题目规定载货、空载里程薪酬。

\subsection{符号说明}
本文统一单位约定：距离单位为 km；速度单位为 km/h；重量单位为 kg；时间统一换算为小时；网点停留时间给定单位为 min，计算时除以 60 转换为小时。
\par 本文中符号说明的表格如下图所示：
\begin{table}[htbp]
\centering
\caption{模型参数与符号说明}
\begin{tabular}{l l l}
\toprule
符号 & 含义 & 单位 \\
\midrule
$V$ & 配送点集合，$V$，$0$ 代表仓库原点 & -- \\
$w_i$ & 第$i$个配送点快件重量 & $\mathrm{kg}$ \\
$Q-max$ & 单车单次最大载重上限，$Q=25$ & $\mathrm{kg}$ \\
$d_{ij}$ & 配送点$i$到$j$的曼哈顿距离 & $\mathrm{km}$ \\
$t-s$ & 单个网点固定卸货停留时间，$t=10$ & $\mathrm{min}$ \\
$T_1$ & 问题一、二单人最大工作时长，$T_1=6$ & $\mathrm{h}$ \\
$T_2$ & 问题三放宽后单人最大工作时长，$T_2=8$ & $\mathrm{h}$ \\
$v_1$ & 问题一全程行驶速度，$v_1=25$ & $\mathrm{km/h}$ \\
$v_2$ & 问题二载货行驶速度，$v_2=20$ & $\mathrm{km/h}$ \\
$v_3$ & 问题二空载返程行驶速度，$v_3=30$ & $\mathrm{km/h}$ \\
$C_1$ & 载货里程计价单价，$C_1=3$ & 元/($\mathrm{km\cdot kg}$) \\
$C_2$ & 空载返程计价单价，$C_2=2$ & 元/$\mathrm{km}$ \\
$K$ & 启用业务员总数量（决策变量） & -- \\
$F_k$ & 第$k$组配送路线总薪酬成本 & 元 \\
$S_k$ & 第$k$个配送趟次包含的配送点集合 & -- \\
$L_k$ & 单条闭合配送路线总里程 & $\mathrm{km}$ \\
$n_k$ & 单趟配送包含网点数量 & 个 \\
\bottomrule
\end{tabular}
\label{tab:para}
\end{table}

\section{问题一 模型建立与求解}

\subsection{问题一数学模型构建}

\subsubsection{距离与工时计算公式}
曼哈顿距离：
\begin{equation}
d_{ij}=|x_i-x_j|+|y_i-y_j|
\label{eq:energy}
\end{equation}

单组路线总工时 = 行驶总时长 + 全部网点停留总时长，单位统一换算为小时：

\begin{equation}
T_{total}=\frac{D_{route}}{v_1}+\frac{n \cdot t_s}{60}
\label{eq:energy}
\end{equation}
其中\(D_{route}\)为单条闭合路线总里程，n为本组配送点数量。

\subsubsection{约束条件}
1.	载重约束：每组所有配送点重量之和不超过单车载重上限
\(\sum_{i \in S_k} w_i \le Q_{max}, \quad \forall k=1,2,...,K\),
\(S_k\)为第k名业务员分配的配送点集合；

\par2. 工时约束：单组总工作时长不超过 6 小时
\(T_{total} \le T_1\)

\par3. 全覆盖约束：全部 30 个配送点仅分配至一个分组，无遗漏、无重复；
\par4. 闭合路线约束：每组路线起点、终点均为仓库原点 0。

\subsubsection{目标函数}
以启用业务员总数最小为优化目标：
\(\min Z = K\)

\subsection{求解算法：最远点优先贪心分组 + 组内最近邻排序}
\subsubsection{算法整体步骤}
步骤 1：输入 30 个配送点坐标、重量，计算全部点到仓库原点的曼哈顿距离；
\par 步骤 2：初始化未分配配送点集合\(U=V \setminus \{0\}\)，分组编号\(k=1\)；
\par 步骤 3：若U为空，算法终止；否则选取U中距离原点最远的点作为本组起点；
\par 步骤 4：循环选取距离当前路线末端最近、加入后总重量≤25kg 的未分配点加入本组；无满足条件点时，完成当前分组；
\par 步骤 5：计算本组闭合路线总里程、总工时，校验工时≤6h；若超时，可回溯遍历当前组内各点，选择移出后可以同时满足载重与工时约束、且整车装载率损失最小的节点移出本趟；
\par 步骤 6：将本组配送点从U中移除，分组编号\(k=k+1\)，返回步骤 3 循环；
\par 步骤 7：每组内部采用最近邻算法排序配送访问顺序，生成完整闭合路线。
\par 步骤 8：采用 FFD 装箱法将所有配送趟次分配给业务员，累加每名业务员各趟工时，总时长不超过 6h，统计最终业务员总数。即将该问题视为装箱子问题：每一个已经生成的配送趟次视为一个物品，物品重量为该趟配送所需总工时；箱子容量为单人单日最大工时（6h 或 8h）；目标为使用最少数量的箱子（业务员）容纳全部物品；采用 FFD（首次适应递减）算法求解该一维装箱问题。

\subsubsection{单组分组示例（第 1 组）}
1.	初始未分配点中距离原点最远点为 30 号（坐标 28,18），累计重量 4.2kg；
\par2.	依次加入相邻可容纳点 29、28、23、11，累计总重量 24.8kg，达到载重上限；
\par3.	路线闭合：0→30→29→28→23→11→0；
\par4.	核算：总里程 106km，总工时 304.4 分钟≤360 分钟（6h），满足全部约束。

\subsection{模型求解过程}
采用 Python 模块化编程完成求解，\(data_points.py\)存储 30 个配送点原始坐标与重量数据；utils.py封装曼哈顿距离、工时核算通用函数；optimizer.py实现贪心分趟逻辑；problem1.py执行完整求解流程；\(plot_vrp.py\)输出配送路线可视化图像；运行后自动导出\(results.csv,results_summary.md\)存储全部调度指标，所有分组逐行校验载重、工时约束，剔除超载、超时无效方案。

\subsection{问题一求解结果}
全局汇总指标：共划分 8 个独立配送趟次，最少配置 5 名业务员，30 点总载货重量 184.5kg，全部配送总行驶里程 478.0km。

\begin{table}[htbp]
\centering
\setlength{\tabcolsep}{4pt} % 缩小列之间空白，核心
\footnotesize % 整体字号变小
\caption{6小时工时约束下各业务员配送调度明细}
\begin{tabularx}{\textwidth}{c c >{\raggedright\arraybackslash}X c c c}
\toprule
业务员编号 & 当日配送趟数 & 覆盖配送点集合 & 当日总载货(kg) & 单日总工时(h) & 单业务员总里程(km) \\
\midrule
1 & 1 & $[30,29,28,23,14]$ & 24.5 & 4.83 & 100.0 \\
2 & 2 & $[26,27,6,11,12,3]$ & 47.8 & 5.88 & 122.0 \\
3 & 2 & $[21,22,10,9,15,8,13,7,5]$ & 44.1 & 5.82 & 108.0 \\
4 & 2 & $[24,25,19,4,2,1]$ & 46.7 & 4.60 & 90.0 \\
5 & 1 & $[18,16,20,17]$ & 21.4 & 2.99 & 58.0 \\
\bottomrule
\end{tabularx}
\label{tab:6h_schedule}
\end{table}


\subsection{结果分析}
1.	分组特征：重量偏大的配送点集中形成独立分组，远距离点位单独分配，避免长距离路线叠加超重、超时；
\par2.	算法优缺点：贪心分组快速得到最小人力可行解，但未考虑配送成本，仅适用于人员成本为核心优化目标的场景；
\par3.	约束校验：全部组别总重量≤25kg、总工时≤360min，模型约束全部满足，方案具备现实可行性。


\section{问题二 模型建立与求解}
\subsection{ 问题二数学模型构建}
\subsubsection{分段薪酬成本计算公式}
设路线访问顺序：\(0 \to p_1 \to p_2 \to ... \to p_n \to 0\),前n段为载货路段，最后一段\(p_n\to0\)为空载路段；分段速度对应的工时计算公式如下：

\begin{equation}
T_{total}=\sum_{i=0}^{n-1}\frac{d_{p_ip_{i+1}}}{v_{load}}+\frac{d_{p_n0}}{v_{empty}}+\frac{n\cdot t_s}{60}
\label{eq:energy}
\end{equation}


\par1.	 载货段成本：每段\(p_i \to p_{i+1}\)行驶时，车上剩余总重量为\(\sum_{t=i+1}^n w_{p_t}\)，载货总费用：
\begin{equation}
Cost_{load}=c_{load} \cdot \sum_{i=0}^{n-1} d_{p_i p_{i+1}} \cdot (\sum_{t=i+1}^n w_{p_t})
\label{eq:energy}
\end{equation}

\par2.	 空载返程成本：最后一点返回仓库里程\(d_{p_n 0}\)，空载费用：
\begin{equation}
Cost_{empty}=c_{empty} \cdot d_{p_n 0}
\label{eq:energy}
\end{equation}

\par3.	 单组总薪酬：
\begin{equation}
Cost = Cost_{load} + Cost_{empty}
\label{eq:energy}
\end{equation}

\subsubsection{约束条件}
载重上限 25kg、单日工时 6h、全覆盖、闭合路线约束与问题一完全一致；仅更新行驶速度参数：载货 20km/h、空载 30km/h，工时计算公式同步更新。

\subsubsection{目标函数}
全部业务员配送薪酬总和最小：
\begin{equation}
\min Z_{cost} = \sum_{k=1}^K Cost_k
\label{eq:energy}
\end{equation}
其中\(Cost_k\)为第k组配送路线薪酬成本。


\subsection{求解算法：基线核算 + 跨组调点局部搜索优化}
\subsubsection{阶段 1：原始分组基准费用计算}
直接复用问题一的 8 组配送路线，按照分段成本公式逐组核算载货、空载费用，汇总得到总薪酬，作为优化对比基线。

\subsubsection{阶段 2：局部搜索优化策略}
在不破坏载重、工时硬约束的前提下，两类邻域操作降低总成本：
\begin{itemize}
    \item MOVE 跨组移动：将某配送点由原组移入另一组，若移动后两组均满足载重、工时约束且总费用下降则接受；
    \item SWAP 跨组交换：交换两组中各一个配送点，在约束可行且总费用下降条件下接受，实现轻重货物在长短途路线间的再分配；
    \item 组内 2-opt 排序：各组访问顺序按曼哈顿距离以最近邻构造、2-opt 优化确定，保证组内行驶里程最短；
    \item 迭代终止条件：连续多轮邻域搜索无成本下降，输出当前最优分组方案。
\end{itemize}

\subsection{求解结果与明细}
全局最优方案汇总：共划分 9 个配送趟次，需配置 6 名业务员，全局最低配送总薪酬 13800.7 元。
\begin{table}[htbp]
\centering
\setlength{\tabcolsep}{4pt}
\footnotesize
\caption{成本最优方案各业务员薪酬明细}
\begin{tabularx}{\textwidth}{c c >{\RaggedRight\arraybackslash}X c c c}
\toprule
业务员编号 & 当日配送趟数 & 覆盖配送点集合 & 当日总载货(kg) & 单日工时(h) & 本组薪酬(元) \\
\midrule
1 & 1 & $[30,29,20,6]$ & 19.9 & 4.50 & 2029.7 \\
2 & 1 & $[26,14,28,5]$ & 24.3 & 4.33 & 2429.8 \\
3 & 2 & $[23,15,13,7,4,8,12,3]$ & 45.3 & 6.00 & 2576.8 \\
4 & 2 & $[24,25,19,18,9]$ & 33.9 & 6.00 & 3046.6 \\
5 & 2 & $[1,11,27,16,17,2]$ & 41.6 & 5.83 & 2460.2 \\
6 & 1 & $[21,22,10]$ & 19.5 & 2.75 & 1257.6 \\
\bottomrule
\end{tabularx}
\label{tab:salary_detail}
\end{table}



\subsection{结果对比与分析}
1.	成本优化逻辑验证：优化方案通过 MOVE/SWAP 跨组调点重组配送分组，并在组内按曼哈顿最短路排序访问顺序，缩短高载重路段的行驶里程；即使优化后总行驶里程小幅上升，载货薪酬下降幅度仍远超空载里程带来的成本增量，最终全局薪酬实现最优；
\par2.	目标冲突分析：最小人力与最小配送成本两个单目标得到的最优调度方案并不相同，仅通过两个独立单目标求解，尚不足以严格证明目标之间存在冲突，可作为拓展研究构建加权多目标优化模型，生成帕累托最优解集，供企业根据经营偏好进行方案权衡；
\par3.	约束校验：优化后所有配送趟次载货≤25kg，每名业务员单日累计工时不超过 6 小时，所有硬约束全部满足，方案合规可行。

\section{问题三 模型建立与求解}
\subsection{模型调整说明}
完整复用问题一整套数学模型、约束、贪心分趟 + FFD 打包算法，仅修改工时约束参数：单人单日最大累计工时由\(T_1=6h\)调整为T\( T_2=8h\)；其余载重、距离、速度、配送点数据全部不变，仅求解最小人力调度方案，与 6 小时工时方案横向对比。
\par 更新后工时约束：
\(∑T_total^((m))≤T_2=8\)

\subsection{求解流程}
1.	仅修改\(config.py\)内工时上限参数为 8，其余代码完全复用\(problem1.py\)；
\par 2.	执行贪心分趟流程，配送趟次划分结果与 6 小时场景一致，仍为 8 趟；
\par 3.	采用 FFD 装箱法重新分配配送趟次给业务员，累加每名业务员各趟总工时，不得超过 8h；
\par 4.	提取业务员总数、总里程、单人平均工时等指标，与 6 小时约束方案对比分析。

\subsection{8 小时工时约束下求解结果}
全局汇总指标：配送趟次保持 8 趟不变，最少仅需 4 名业务员完成全部配送，全局总行驶里程维持 478.0km 不变。

\begin{table}[htbp]
\centering
\setlength{\tabcolsep}{4pt}
\footnotesize
\caption{8小时工时约束下各业务员调度明细}
\begin{tabularx}{\textwidth}{c c >{\RaggedRight\arraybackslash}X c c}
\toprule
业务员编号 & 当日配送趟数 & 覆盖配送点集合 & 单日总工时(h) & 单业务员总里程(km) \\
\midrule
1 & 2 & $[30,29,28,23,14,18,16,20,17]$ & 7.82 & $158$ \\
2 & 2 & $[26,27,6,21,22,10,9,15]$ & 6.85 & $138$ \\
3 & 3 & $[24,25,19,8,13,7,5,4,2,1]$ & 7.11 & $136$ \\
4 & 1 & $[11,12,3]$ & 2.34 & $46$ \\
\bottomrule
\end{tabularx}
\label{tab:schedule_8h}
\end{table}

\subsection{6h 与 8h 工时约束多指标对比分析}
\begin{table}[htbp]
\centering
\setlength{\tabcolsep}{4pt}
\footnotesize
\caption{不同工时约束下最小人力方案指标对比}
\begin{tabularx}{\textwidth}{l l l >{\RaggedRight\arraybackslash}X}
\toprule
评价指标 & 单人工时上限6h & 单人工时上限8h & 指标变化 \\
\midrule
所需最少业务员数量 & 5人 & 4人 & 减少1人，人力投入下降20\% \\
配送总趟次 & 8趟 & 8趟 & 无变化，载重约束决定趟次下限 \\
全局总行驶里程 & 478.0\,km & 478.0\,km & 无变化，由配送点空间分布决定 \\
单人最大单日工时 & 5.88\,h & 7.82\,h & 单人工作负荷显著提升 \\
\bottomrule
\end{tabularx}
\label{tab:compare_6h_8h}
\end{table}

\subsection{结果综合分析}
1.	人力优化规律：工时上限属于强有效约束，放宽单人每日允许工作时长，单名业务员可承接更多配送趟次，直接降低企业所需业务员数量，减少固定人力成本；
\par 2.	里程不变：在本算例保持原有配送趟次分组不变的前提下，总行驶里程维持 478.0 km 不变；
\par 3.	运营权衡建议：延长工作时长可节约人力，但会提升业务员单日配送负荷，存在疲劳作业风险；企业需要结合劳动法规、加班薪资、人员流失成本综合确定最优单日工时标准，不可单纯依靠延长工时压缩用工人数。


\section{模型检验与有效性分析}
\subsection{算法约束校验检验}
对问题一、问题二全部分组逐条核验两类硬约束：
\begin{itemize}
    \item 载重校验：每组总重量均≤25kg，无超载分组；
    \item 工时校验：6 小时约束下所有路线总工时均≤360 分钟，无超时排班；
    \item 全覆盖校验：30 个配送点全部分配，无遗漏、无重复分配；
    \item 闭合路线校验：所有路线起点、终点均为仓库原点，满足往返配送要求。
\end{itemize}
	全部方案均满足题目硬性规则，算法输出解具备可行性。

\subsection{小规模手工验算检验}
选取前 5 个配送点构建小规模测试算例，手工计算曼哈顿距离、单趟工时、分段薪酬成本，计算结果与程序输出完全一致，验证距离、工时、成本计算公式代码逻辑无错误。

\subsection{算法稳定性说明}
确定性算法，结果恒一致，每次运行\(problem1.py,problem2.py,problem3.py\)输出的配送分组、路线、成本、人员数量完全一致，模型算法具备稳定可复现性，实验结果可信。

\section{模型评价与改进方向}
\subsection{模型优点}
1.	分层递进建模框架通用性强，三问共用底层计算模块，仅调整目标函数、工时参数即可完成多场景求解，代码复用度高；
\par2.	算法选择由简到繁，优先采用贪心启发式算法，代码易实现、计算速度快，适配 30 点中小规模配送场景；
\par3.	完整覆盖载重、工时双重硬性运营约束，所有输出调度方案均满足题目全部规则，贴合快递末端配送真实运营场景；
\par4.	第三问开展工时参数灵敏度分析，量化工时约束对人力配置的影响，为企业排班管理提供量化决策支撑。

\subsection{模型缺点}
1.	贪心算法、2-opt 邻域搜索/跨路径 swap 操作局部搜索均为启发式算法，仅能输出较优可行解，无法严格证明全局最优；
\par 2.	模型为静态调度模型，假设所有配送订单提前完整已知，未考虑动态临时订单、网点配送时间窗约束；
\par 3.	成本体系简化，仅计算里程相关薪酬，未纳入业务员底薪、车辆购置、燃油损耗、加班补贴等完整企业运营成本；
\par 4.	未区分业务员工作负荷均衡性，8 小时场景存在业务员单日配送趟次差异较大的不均衡问题。

\subsection{模型改进方向}
1.	高精度寻优：以贪心算法输出方案为初始解，引入遗传算法、模拟退火元启发式算法全局迭代，逼近最小人力、最小成本全局最优解；
\par 2.	场景拓展优化：新增网点时间窗、动态临时订单、多车型差异化载重约束，贴合城市快递真实配送场景；
\par 3.	成本体系完善：加入车辆固定出车成本、业务员基础薪资、加班溢价，构建全维度企业综合运营成本模型；
\par 4.	精确求解方案：可构建 0‑1 整数线性规划模型，但该模型仅适合 10 个节点以内小规模场景使用 Gurobi、CPLEX 求取精确最优解；对于本题 30 个节点规模，精确求解难度极大，工程上一般采用元启发式算法求取高质量可行解；
\par 5.	均衡优化拓展：新增业务员单日工时均衡目标，在最小人力 / 最小成本基础上，平衡各业务员工作负荷。




%参考文献
\begin{thebibliography}{9}%宽度9
    \bibitem[1]{liuhaiyang2013latex}
    谢金星，薛毅.优化建模与 LINDO/LINGO 软件 [M], 北京,清华大学出版社，2005.
    \bibitem[2]{mathematical-modeling}
    刘长石，陈慧璇.城市物流配送的混合车辆路径规划模型与优化算法 [J],控制与决策，2023.
     \bibitem[3]{liuhaiyang2013latex}
    Clarke G, Wright J W. Scheduling of vehicles from a central depot to a number of delivery points [J]. Operations Research,1964.
    \bibitem[4]{mathematical-modeling}
    祁文祥, 陆志强, 孙小明. 带软时间窗的集货与送货多车辆路径问题节约算法[J]. 交通运输工程学报, 2010, 10(2): 99-103. doi: 10.19818/j.cnki.1671-1637.2010.02.018
\end{thebibliography}

\newpage
%附录
\begin{appendices}

\section{30 个配送点原始坐标与单件重量数据表}
\begin{figure}[!h]
    \centering
    \includegraphics[width=0.5\linewidth]{屏幕截图 2026-08-13 183545.png}
    \label{fig:placeholder}
\end{figure}

\begin{figure}[!h]
    \centering
    \includegraphics[width=0.5\linewidth]{屏幕截图 2026-08-13 183552.png}
    \label{fig:placeholder}
\end{figure}

\section{全套 Python 求解代码}
\begin{lstlisting}[language=matlab]
1.	data_points.py：30 点原始数据定义与总重量自检；
2.	config.py：全局参数配置（载重、工时、速度、计价单价）；
3.	utils.py：曼哈顿距离、工时、分段薪酬通用计算函数；
4.	optimizer.py：最远点优先贪心分趟、MOVE/SWAP 局部搜索模块；
5.	problem1.py：6 小时最小人力调度求解程序；
6.	problem2.py：分段计价最小薪酬优化求解程序；
7.	problem3.py：8 小时工时约束调度对比求解程序；
8.	plot_vrp.py：配送路线可视化绘图代码。
 \end{lstlisting}

 \section{中间计算数据文件}
\begin{lstlisting}[language=c]
1.	results.csv：三问全部业务员、趟次、重量、工时、里程、薪酬汇总表；
2.	results_summary.md：标准化方案文字汇总；
3.	30×30 配送点曼哈顿距离矩阵；
4.	各配送趟次网点访问路线完整明细。
\end{lstlisting}

 \section{中间计算数据文件}
\begin{figure}[!h]
    \centering
    \includegraphics[width=0.5\linewidth]{q1_strategy.png}
    \caption{问题一 6 小时工时 5 名业务员配送路线分布图}
    \label{fig:placeholder}
\end{figure}

\begin{figure}[!h]
    \centering
    \includegraphics[width=0.5\linewidth]{q2_cost.png}
    \caption{问题二成本最优 6 名业务员配送路线分布图}
    \label{fig:placeholder}
\end{figure}

\begin{figure}
    \centering
    \includegraphics[width=0.5\linewidth]{q3_8h.png}
    \caption{问题三 8 小时工时 4 名业务员配送路线分布图}
    \label{fig:placeholder}
\end{figure}


\end{appendices}

\end{document} 